\documentclass[11pt]{article}
\usepackage[utf8]{inputenc}
\usepackage[polish]{babel}
\usepackage[QX]{fontenc}
\usepackage{lmodern}
\usepackage{textcomp}
\usepackage{graphicx}

 

\title{}
\author{
  Daniel Danilewicz \\
  Maciej Kucharski \\
  Marek Dzienisiuk \\
  Paweł Stiasny (kierownik) \\
  Piotr Dziewicki \\
  Tomasz Pęksa \\
  Wiktor Garbacz \\
}


\begin{document}

\begin{titlepage}
  \centering
  \includegraphics[scale=0.3]{logopw.pdf} \par
  \vspace{0.5cm}

  {\Large Projekt RSO \par 2016L \par}
  \vspace{1cm}
  {\LARGE ,,Usługa bezpiecznej niezawodnej dystrybucji przetworzonej chronionej
          informacji'' \par}

  \vspace{2cm}

  Daniel Danilewicz \par
  Maciej Kucharski \par
  Marek Dzienisiuk \par
  Paweł Stiasny \par
  Piotr Dziewicki \par
  Tomasz Pęksa \par
  Wiktor Garbacz \par

  \vfill
\end{titlepage}

\section{Wymagania}

Usługa powinna:

\begin{tabular}{ l p{10cm} }
  \hline
  RQ1 & Mieć dokładnie jeden ten sam tekstowy plik konfiguracyjny dla części
    serwerowej i klientów usługi, \\ \hline
  RQ2 & Składać się po stronie serwerowej z dwóch warstw: (1)  wewnętrznej 
    wytwarzającej, przechowującej i przesyłającej przetworzoną wrażliwą
    informację, (2) warstwy zewnętrznej do udostępniania przetworzonej informacji
    oprogramowaniu klienckiemu usługi, \\ \hline
  RQ3 & Oprogramowanie realizujące część serwerową powinno być współbieżne, \\
    \hline
  RQ4 & Składać się po stronie klienckiej z prostych narzędzi opatulających
    zaproponowane API protokołu komunikacyjnego, \\ \hline
  RQ5 & Zapewniać odporność na uszkodzenia węzła, stąd każda z warstw powinna być
    uruchomiona na co najmniej dwóch węzłach, \\ \hline
  RQ6 & Zapewniać realizację usługi w trakcie awarii w czasie obsługi, \\ \hline
  RQ7 & Obsługiwać scenariusz próby ponownego wpięcia się przez węzeł dowolnej z
    warstw, z którym wcześniej utracono łączność, a zarazem być odporną na próbę
    zastąpienia nieosiągalnego sieciowo węzła (np. w wyniku ataku DoS, odmowa
    usługi) węzłem wrogim do tego nieuprawnionym, \\ \hline
  RQ8 & Na bieżąco i transparentnie dla oprogramowania klienckiego zarządzać
    dostępnymi zasobami lokalizacyjnymi, \\ \hline
  RQ9 & Zapewniać zadany (większy niż 1 ale dopuszczalny konfiguracją mniejszy niż
    liczba serwerów danych) poziom redundancji - z obsługą uzupełniania kopii na
    innych serwerach w przypadku awarii włącznie, \\ \hline
  RQ10 & Cała część sewerowa usługi powinna być startowana i zamykana jednokrotnym
    wywołaniem skyptu na jednym węźle serwera z jednym argumentem
    \textlangle{} start \textbar{} stop \textbar{} status \textrangle{}
    (wykorzystanie nieinteraktywnych metod automatycznego uwierzytelniania
    węzłów w SSH), \\ \hline
\end{tabular}
 
\end{document}
